\documentclass[../../../include/open-logic-section]{subfiles}

\begin{document}

\olfileid{sth}{story}{approach}
\olsection{في النهج التراكمي-التكراري}

وفيما يأتي بيان أتم قليلًا لكيفية ترتيب المجموعات في مراحل:
\begin{quote}
	تتكون المجموعات في \emph{مراحل}. ولكل مرحلة~${\OLId{latin-looped}{S}}$ مراحل معينة
	\emph{تسبق}~${\OLId{latin-looped}{S}}$. وفي المرحلة~${\OLId{latin-looped}{S}}$، يتحول كل تجمع مؤلف من مجموعات
	تكونت في مراحل سابقة لـ${\OLId{latin-looped}{S}}$ إلى مجموعة. ولا توجد مجموعات غير
	المجموعات التي تتكون في مراحل. \citep[p.~323]{Shoenfield:AST}
\end{quote}
وهذا هو الرسم العام لـ\emph{التصور التراكمي-التكراري للمجموعة}، ومنه
نستمد أصول نظرية المجموعات الصورية المبسوطة في
\olref[sth][][]{part}.

ولنفصل الأمر. بحسب وصف شونفيلد نكوّن في كل مرحلة مجموعات جديدة من
المجموعات التي سبقت إتاحتها. فالمرحلة الأولى، أي المرحلة~${\OLMathNumeral{0}}$، لا
تسبقها في صورته مرحلة، فلا تتاح فيها مجموعة من مرحلة سالفة.
\footnote{لماذا ينبغي أن نفترض أن مرحلة أولى
  \emph{توجد}؟ انظر حاشية~\stagesord{} في~\olref[z][story]{sec}.}
فلا نكوّن إلا مجموعة واحدة لا !!{element} فيها، وهي~$\emptyset$. وفي
المرحلة~${\OLMathNumeral{1}}$ لا تتاح من المراحل السابقة إلا مجموعة واحدة، فتكون المجموعة
الجديدة الوحيدة~$\{\emptyset\}$. وفي المرحلة~${\OLMathNumeral{2}}$ تتاح مجموعتان، فنكوّن
الجديدتين~$\{\{\emptyset\}\}$ و$\{\emptyset, \{\emptyset\}\}$. وعند
المرحلة~${\OLMathNumeral{3}}$ تتاح أربع مجموعات، فنكوّن اثنتي عشرة مجموعة جديدة\ldots.
وعلى هذا تقارب صورة المجموعات التراكمية-التكرارية الرسم الآتي، والأعداد
فيه أسماء المراحل:
\begin{center}
	\begin{tikzpicture}[scale=0.6]
	\tikzset{cut_here/.style={densely dotted}}
	\draw (-4,6) -- (-1, 0)-- (2,6); 
	\draw[cut_here] (-5,8)--(-4,6);
	\draw[cut_here] (2,6)--(3,8);
	\draw(-1.5, 1)--(-0.5, 1);
	\draw(-2, 2)--(0, 2);
	\draw(-2.5, 3)--(0.5,3);
	\draw(-3, 4)--(1, 4);
	\draw(-3.5, 5)--(1.5, 5);
	\draw(-4, 6)--(2, 6);
	\node[label] at (0, 0) {\small 0};
	\node[label] at (0.5, 1) {\small 1};
	\node[label] at (1, 2) {\small 2};
	\node[label] at (1.5, 3) {\small 3};
	\node[label] at (2, 4) {\small 4};
	\node[label] at (2.5, 5) {\small 5};
	\node[label] at (3, 6) {\small 6};
	\end{tikzpicture}
\end{center}
فما الذي يدعونا إلى قبول هذه الرواية؟

من أسباب ذلك أنها صورة حسنة سهلة العمل؛ وبعد بطلان الصورة الأظهر، أعني
الاستيعاب الساذج، لا غنى لنا عن بديل صالح.

وليس حسنها وحده حجة لها، بل ثمة سبب قوي للظن بأن كل نظرية مجموعات تُبنى
عليها تكون \emph{متسقة}. وبيانه:

نكوّن المجموعات، في التصور التراكمي-التكراري، على مراحل، ولا تكون
!!{element} إلا أشياء أتيحت \emph{من قبل}. ولذلك نستطيع في كل
مرحلة~${\OLId{latin-looped}{S}}$ أن نكوّن المجموعة
\[
	{\OLId{latin-looped}{R}}_{\OLId{latin-looped}{S}} = \Setabs{{\OLId{latin-ordinary}{x}}}{{\OLId{latin-ordinary}{x}} \notin {\OLId{latin-ordinary}{x}} \text{ و${\OLId{latin-ordinary}{x}}$ كان متاحًا قبل ${\OLId{latin-looped}{S}}$}}
\]
فيبرهن استدلال مفارقة راسل أن~${\OLId{latin-looped}{R}}_{\OLId{latin-looped}{S}}$ ذاتها لم تكن متاحة قبل
المرحلة~${\OLId{latin-looped}{S}}$، ولا تناقض في ذلك. بل إن قبول التصور التراكمي-التكراري يمنع
أن \emph{نتوقع} إنشاء مجموعة راسل أصلًا؛ لأنها كانت ستجمع كل مجموعة لا
تنتمي إلى نفسها «مما سيتاح في وقت ما». فثبوت عجزنا عن إنشاء مجموعة راسل
ليس \emph{مستغربًا} في هذه الرواية، وإنما هو عين ما \emph{تتنبأ} به.

%من جهة ما، يقدم التصور التراكمي-التكراري للمجموعة جوابًا عن مفارقة راسل يشبه جواب \emph{أنصار الحملية}. فقد قال أنصار الحملية إن تجمع جميع المجموعات$_n$ التي لا تنتمي إلى نفسها هو مجموعة$_{n+1}$، وإن هذه المجموعة$_n$ لن تنتمي إلى نفسها. (بل يعد أكثر أنصار الحملية محاولة السؤال عما إذا كانت مجموعة تنتمي إلى نفسها \emph{غير سليمة نحويًا}.)
%
%لكن هناك فرقًا مهمًا بين النهج التراكمي-التكراري والنهج الحملي: يعامل النهج التراكمي-التكراري جميع كياناتنا بوصفها \emph{من النوع نفسه}. ولدى أنصار الحملية، على الأرجح، مجموعة خالية$_0$ (أي مجموعة$_0$ لا !!{element} فيها)، ومجموعة خالية$_1$ (أي مجموعة$_1$ لا !!{element} فيها)، وهما \emph{كيانان مختلفان}. أما في النهج التراكمي-التكراري فليست هناك إلا مجموعة خالية واحدة هي~$\emptyset$، وستكون متاحة في كل مرحلة من الهرم.

%وبالفعل، ستصدق مسلّمة الامتدادية، التي قُررت في بداية هذا الفصل، على عناصر هذا الهرم (بحيث لا يمكن أن توجد إلا مجموعة «خالية» \emph{واحدة}). لكنني لا أنوي استعمال التصور التراكمي-التكراري لتبرير الامتدادية (انظر~\cite{Boolos1971}). وأكرر أنني أقترح قبول الامتدادية لمجرد اهتمامنا بالتجمعات الامتدادية.

\end{document}



